% GENERATED FILE. Edit canonical lessons, then run npm run generate:books.
% canonical-source-hash: fnv1a64:967a1fa2f4512ead
% canonical-lessons: FR-C329-saladier, FR-C329-verser, FR-C329-battre, FR-C329-fouet, FR-C329-fouetter

\chapter{Battre les œufs}
\label{ch:fr-battre-les-ufs}
\hlchaptermodality{329}

\begin{chapteropening}
\textbf{By the end of this chapter:} \emph{I can describe mixing: the mixing bowl, pouring, beating, the whisk and whisking.}

Complete the last lesson of chapter 329: I can describe mixing: the mixing bowl, pouring, beating, the whisk and whisking.
\end{chapteropening}

\section[le saladier]{le saladier — the salad bowl, the mixing bowl}

\label{lesson:FR-C329-saladier}



\begin{quote}
Before the new one: say the French for to slice thinly, then the French for to grate.
\end{quote}

\subsection*{You'll want to know: le saladier}
\textbf{le saladier} — \textquotedblleft{}the salad bowl, the mixing bowl\textquotedblright{}.

\begin{quote}
Elle met la farine dans un grand saladier.
\end{quote}

\begin{quote}
\emph{She puts the flour in a big bowl.}
\end{quote}

Built on \emph{salade}.

\begin{cousinweb}
\emph{saladier} is built on \emph{salade}, which EtymDB traces to Italian \emph{salata} \textquotedblleft{}something salted\textquotedblright{} (Provençal \emph{salada} is also proposed), built on Latin \emph{sal} \textquotedblleft{}salt\textquotedblright{} (Lewis \& Short).
\end{cousinweb}

\begin{grammarlens}[title={What you've built}]
One more word for telling, explaining and describing at B1.
\end{grammarlens}

\subsection*{Your turn}
\begin{itemize}
\raggedright
  \item \emph{Say it:} \emph{le saladier}
  \item \emph{Say it:} \emph{le saladier}, once more
  \item \emph{Say it:} a sentence of your own with \emph{le saladier}
  \item \emph{Recall:} say the French for to slice thinly, then the French for to grate, then say \emph{le saladier} again
\end{itemize}

\begin{tcolorbox}[breakable,colback=teal!4,colframe=teal!35!black,title={Before you move on}]
What does le saladier mean? (\textquotedblleft{}The salad bowl, the mixing bowl\textquotedblright{}.) Say it once more.
\end{tcolorbox}
\section[verser]{verser — to pour; to pay in}

\label{lesson:FR-C329-verser}



\begin{quote}
Before the new one: say the French for to grate, then the French for the salad bowl.
\end{quote}

\subsection*{You'll want to know: verser}
\textbf{verser} — \textquotedblleft{}to pour; to pay in\textquotedblright{}.

\begin{quote}
Il verse le lait dans le saladier.
\end{quote}

\begin{quote}
\emph{He pours the milk into the bowl.}
\end{quote}

Also to pay money in: \emph{verser un salaire}.

\begin{cousinweb}
\emph{verser} is inherited from Latin \emph{versāre} \textquotedblleft{}to turn often\textquotedblright{} (EtymDB), built on \emph{vertere} \textquotedblleft{}to turn\textquotedblright{}.
\end{cousinweb}

\begin{grammarlens}[title={What you've built}]
One more word for telling, explaining and describing at B1.
\end{grammarlens}

\subsection*{Your turn}
\begin{itemize}
\raggedright
  \item \emph{Say it:} \emph{verser}
  \item \emph{Say it:} \emph{verser}, once more
  \item \emph{Say it:} a sentence of your own with \emph{verser}
  \item \emph{Recall:} say the French for to grate, then the French for the salad bowl, then say \emph{verser} again
\end{itemize}

\begin{tcolorbox}[breakable,colback=teal!4,colframe=teal!35!black,title={Before you move on}]
What does verser mean? (\textquotedblleft{}To pour; to pay in\textquotedblright{}.) Say it once more.
\end{tcolorbox}
\section[battre]{battre — to beat}

\label{lesson:FR-C329-battre}



\begin{quote}
Before the new one: say the French for the salad bowl, then the French for to pour.
\end{quote}

\subsection*{You'll want to know: battre}
\textbf{battre} — \textquotedblleft{}to beat\textquotedblright{}.

\begin{quote}
Elle bat les œufs avec une fourchette.
\end{quote}

\begin{quote}
\emph{She beats the eggs with a fork.}
\end{quote}

\emph{Battre les blancs en neige} — to beat egg whites until stiff. Also to beat someone at a game.

\begin{cousinweb}
\emph{battre} is inherited from Late Latin \emph{battere} \textquotedblleft{}to beat\textquotedblright{} (EtymDB), from Latin \emph{battuere} (Lewis \& Short).
\end{cousinweb}

\begin{grammarlens}[title={What you've built}]
One more word for telling, explaining and describing at B1.
\end{grammarlens}

\subsection*{Your turn}
\begin{itemize}
\raggedright
  \item \emph{Say it:} \emph{battre}
  \item \emph{Say it:} \emph{battre}, once more
  \item \emph{Say it:} a sentence of your own with \emph{battre}
  \item \emph{Recall:} say the French for the salad bowl, then the French for to pour, then say \emph{battre} again
\end{itemize}

\begin{tcolorbox}[breakable,colback=teal!4,colframe=teal!35!black,title={Before you move on}]
What does battre mean? (\textquotedblleft{}To beat\textquotedblright{}.) Say it once more.
\end{tcolorbox}
\section[le fouet]{le fouet — the whisk; the whip}

\label{lesson:FR-C329-fouet}



\begin{quote}
Before the new one: say the French for to pour, then the French for to beat.
\end{quote}

\subsection*{You'll want to know: le fouet}
\textbf{le fouet} — \textquotedblleft{}the whisk; the whip\textquotedblright{}.

\begin{quote}
Il mélange la crème avec un fouet.
\end{quote}

\begin{quote}
\emph{He mixes the cream with a whisk.}
\end{quote}

\emph{Un coup de fouet} — a boost.

\begin{cousinweb}
\emph{fouet} is usually explained as a diminutive of Old French \emph{fou} \textquotedblleft{}beech\textquotedblright{}, from Latin \emph{fāgus} \textquotedblleft{}beech\textquotedblright{} (Lewis \& Short): first a beech switch.
\end{cousinweb}

\begin{grammarlens}[title={What you've built}]
One more word for telling, explaining and describing at B1.
\end{grammarlens}

\subsection*{Your turn}
\begin{itemize}
\raggedright
  \item \emph{Say it:} \emph{le fouet}
  \item \emph{Say it:} \emph{le fouet}, once more
  \item \emph{Say it:} a sentence of your own with \emph{le fouet}
  \item \emph{Recall:} say the French for to pour, then the French for to beat, then say \emph{le fouet} again
\end{itemize}

\begin{tcolorbox}[breakable,colback=teal!4,colframe=teal!35!black,title={Before you move on}]
What does le fouet mean? (\textquotedblleft{}The whisk; the whip\textquotedblright{}.) Say it once more.
\end{tcolorbox}
\section[fouetter]{fouetter — to whisk, to whip}

\label{lesson:FR-C329-fouetter}



\begin{quote}
Before the new one: say the French for to beat, then the French for the whisk.
\end{quote}

\subsection*{You'll want to know: fouetter}
\textbf{fouetter} — \textquotedblleft{}to whisk, to whip\textquotedblright{}.

\begin{quote}
Elle fouette la crème pendant cinq minutes.
\end{quote}

\begin{quote}
\emph{She whips the cream for five minutes.}
\end{quote}

\emph{La crème fouettée} — whipped cream.

\begin{cousinweb}
\emph{fouetter} is built on \emph{fouet} \textquotedblleft{}whip, whisk\textquotedblright{} (EtymDB).
\end{cousinweb}

\begin{grammarlens}[title={What you've built}]
One more word for telling, explaining and describing at B1.
\end{grammarlens}

\subsection*{Your turn}
\begin{itemize}
\raggedright
  \item \emph{Say it:} \emph{fouetter}
  \item \emph{Say it:} \emph{fouetter}, once more
  \item \emph{Say it:} a sentence of your own with \emph{fouetter}
  \item \emph{Recall:} say the French for to beat, then the French for the whisk, then say \emph{fouetter} again
\end{itemize}

\begin{tcolorbox}[breakable,colback=teal!4,colframe=teal!35!black,title={Before you move on}]
What does fouetter mean? (\textquotedblleft{}To whisk, to whip\textquotedblright{}.) Say it once more.
\end{tcolorbox}
